\documentclass[UTF8,11pt,oneside]{ctexbook}
\usepackage[a4paper,margin=2.35cm,headheight=15pt]{geometry}
\usepackage{amsmath,amssymb,mathtools,booktabs,longtable,array,multicol}
\usepackage{xcolor,graphicx,enumitem,fancyhdr,listings,hyperref}
\definecolor{navy}{HTML}{17324D}
\definecolor{teal}{HTML}{197A78}
\definecolor{pale}{HTML}{F1F6F7}
\definecolor{codebg}{HTML}{F5F7F9}
\hypersetup{colorlinks=true,linkcolor=navy,urlcolor=teal,pdftitle={机器学习零基础讲义：导论与回归},pdfauthor={个人学习讲义}}
\ctexset{chapter={format=\Huge\bfseries\color{navy}},section={format=\Large\bfseries\color{navy}},subsection={format=\large\bfseries\color{teal}}}
\pagestyle{fancy}\fancyhf{}\fancyhead[L]{机器学习零基础讲义}\fancyhead[R]{导论与回归}\fancyfoot[C]{\thepage}
\setlist[itemize]{leftmargin=2em,itemsep=2pt,topsep=3pt}
\setlist[enumerate]{leftmargin=2em,itemsep=2pt,topsep=3pt}
\lstset{basicstyle=\ttfamily\small,backgroundcolor=\color{codebg},frame=single,breaklines=true,showstringspaces=false,columns=fullflexible,language=Python,keywordstyle=\color{navy}\bfseries,commentstyle=\color{teal},numbers=left,numberstyle=\tiny\color{gray},xleftmargin=1.7em,framexleftmargin=1.7em}
\newcommand{\key}[1]{\textcolor{teal}{\bfseries #1}}
\newcommand{\R}{\mathbb{R}}
\newcommand{\E}{\mathbb{E}}
\newcommand{\Var}{\operatorname{Var}}
\newcommand{\MSE}{\operatorname{MSE}}
\newcommand{\RMSE}{\operatorname{RMSE}}
\newcommand{\MAE}{\operatorname{MAE}}
\newenvironment{remember}{\begin{quote}\color{teal}\textbf{记住这一点：}}{\end{quote}}
\begin{document}
\hypersetup{pageanchor=false}
\begin{titlepage}
\centering\vspace*{4cm}
{\Huge\bfseries\color{navy}机器学习零基础讲义\par}
\vspace{0.8cm}{\LARGE 导论与回归\par}
\vspace{1.6cm}{\large 面向没有机器学习基础、Python 基础一般的自学者\par}
\vfill
{\large 根据《机器学习简介 · 第一讲》《机器学习导论 · 第 2 讲：回归》\par
及《回归：从条件期望到正则化》整理、核算并补充\par}
\vspace{1cm}{\large 个人学习使用 · 2026 年 9 月\par}
\end{titlepage}
\hypersetup{pageanchor=true}
\frontmatter

\chapter*{这份讲义用来做什么}
\addcontentsline{toc}{chapter}{这份讲义用来做什么}
这是一份\key{个人自学和复习讲义}：把两讲课件中的概念、案例和回归方法串成可以独立阅读的路线，并用第三份 PDF 补足推导。它适合课前预习、课后复习、做题时查公式和第一次运行 Python 例子。它覆盖\key{机器学习导论与回归}，并不是整门机器学习课程的全部；分类、聚类、神经网络训练等后续内容在此只建立地图，不作完整课程展开。

\textbf{建议用法：}第一次读，依次读第 1--7 章和第 10 章，先理解句子、表格与手算；第 8--9 章的矩阵推导和优化可以第二遍再读；第 11 章代码边运行边改数字；最后做第 12 章自测。遇到公式先问：\emph{输入是什么、输出是什么、每个符号代表什么、它解决哪个问题}，再看计算。

\textbf{范围与来源：}第一讲约 61 张幻灯片，第二讲约 62 张幻灯片；第三份 21 页 PDF 聚焦回归。本文用自己的例子和表述重新组织内容。课件中的案例和表格只用于教学演示，不能作为真实房价或产品性能证据。文末列出三份本地来源和对应章节。

\tableofcontents
\mainmatter

\chapter{先认识机器学习在做什么}
\section{一个不神秘的定义}
Tom Mitchell 的常用表述是：如果程序在任务 $T$ 上的表现 $P$ 随经验 $E$ 增加而改善，就说它从经验中学习。比如识别垃圾邮件：$T$ 是判断邮件类别，$P$ 可以是误报率，$E$ 是已经标注的邮件。三个字母能防止一句空泛的“模型更聪明了”掩盖评估问题。

传统程序通常由人写明规则；机器学习程序由人指定问题、数据表示、模型范围和评价方式，让算法从样本中估计参数。人仍然要负责\key{收集什么数据、标签是否可靠、犯哪种错更贵、部署时是否还能用}。

\section{统一的五步链条}
给定输入 $x$，模型 $f_\theta$ 产生预测 $\hat y=f_\theta(x)$。将预测与正确答案 $y$ 比较，得到损失 $\ell(y,\hat y)$；优化算法调整参数 $\theta$；最后在未见过的数据上检验。简写为
\[
\boxed{\text{数据}\longrightarrow\text{模型}\longrightarrow\text{损失}\longrightarrow\text{优化}\longrightarrow\text{泛化检验}}.
\]
例如房价预测：面积和房龄是输入，成交价是目标；$\hat y=wx+b$ 是模型；$(y-\hat y)^2$ 是单次平方损失；寻找合适的 $w,b$ 是训练；在新房子上测试是泛化检验。

\section{三种学习范式}
\begin{center}\begin{tabular}{p{2.4cm}p{4.7cm}p{6.1cm}}\toprule
范式&数据与反馈&典型问题\\\midrule
监督学习&每个输入有目标或标签&回归预测连续值；分类预测类别或类别概率\\
无监督学习&输入没有预先给定的标签&聚类寻找分组；PCA 寻找保留较多变化的方向\\
强化学习&行动改变环境，奖励可能延迟&游戏落子、机器人动作、长期决策\\\bottomrule\end{tabular}\end{center}
“没有标签”并不等于“没有目标”：K-means 仍优化簇内距离，PCA 仍优化投影后的方差。强化学习则要考虑当前行动怎样影响未来能看到的数据。

\section{用课件案例读懂共同框架}
\begin{itemize}
\item \textbf{DQN 与 AlphaGo：}游戏画面或棋盘是状态，动作会影响后续状态，胜负或得分给出反馈。AlphaGo 还把学到的策略、价值估计与树搜索组合。AlphaZero 借助可准确模拟的棋类规则进行大量自我对弈；现实道路或机器人通常没有这样低成本、完全可靠的试错环境。
\item \textbf{AlphaFold：}输入蛋白质序列等信息，输出三维结构，利用已知结构训练。局部置信度与实验验证仍重要。
\item \textbf{语言模型：}给前面的 token 预测下一个 token，文本自身提供训练目标，属于自监督学习的一种。语言流畅不保证事实正确。注意力机制使每个位置按当前需要汇总其他位置的信息；先记住这个作用，不必先推导 Transformer。
\item \textbf{图像与视频生成：}根据文字条件生成可能的图像或视频，目标是学习一个条件分布，所以同一句描述可以产生不同结果。视频还需维持跨帧的一致性。
\item \textbf{自动驾驶与机器人：}感知、预测、规划、控制形成闭环。平均准确率无法说明低概率高后果事件的风险；训练环境与真实环境之间可能发生分布偏移。
\end{itemize}
对任何新系统都可以追问：\emph{它看过什么、优化什么、怎样测试、在哪些场景可靠、失败会怎样}。这比背产品名称更耐用。

\section{常见分支怎样放进这张地图}
\textbf{深度学习}让多层函数从数据中学习表示。CNN 复用局部卷积核，常见于图像；RNN 按时间更新隐藏状态，适合序列；Transformer 用注意力汇总不同位置的信息，广泛用于语言和多模态。两层纯线性变换仍能合并为一层，所以神经网络需要 ReLU 等非线性激活函数。模型更大并不自动解决数据质量、标注偏差或评估问题。

\textbf{计算机视觉}的输出形式各不相同：分类给整张图一个类别，检测给位置和类别，分割给每个像素类别。少数类重要时，只看准确率会误导。例如 100 张图仅有 2 张缺陷图，全部预测“正常”仍有 98\% 准确率，却完全找不到缺陷。此时要看精确率 $TP/(TP+FP)$ 与召回率 $TP/(TP+FN)$，其中 $TP$ 是正确报出的正例，$FP$ 是误报，$FN$ 是漏报。

\textbf{自然语言处理}处理文本的上下文与组合含义。检索增强生成（RAG）先从外部文档检索证据，再据此生成答案；若检索遗漏或证据本身错误，生成结果仍可能不可靠。\textbf{具身智能}把视觉、语言和动作放进反馈闭环；VLA（视觉—语言—动作）模型根据观察与指令输出动作，但真实机器人轨迹采集昂贵，仿真到现实还存在差距。这些系统都需要定义任务、数据、目标和失败检查。

\chapter{读公式前的数学与 Python 预备}
\section{变量、函数、求和}
$x$ 是输入，$y$ 是真实答案，$\hat y$ 读作“预测的 $y$”。函数 $f(x)$ 是一个规则：给定 $x$，返回一个值。$\sum_{i=1}^{n}a_i$ 表示把 $a_1$ 到 $a_n$ 相加；$\frac1n\sum_i a_i$ 是平均。$\arg\min_w J(w)$ 表示“使 $J$ 最小的那个 $w$”，其结果是参数，不是最小值本身。

\section{向量和矩阵只需先会看形状}
把一个样本的 $d$ 个特征写成 $x_i=(x_{i1},\ldots,x_{id})^\top\in\R^d$。把 $n$ 个样本按行堆起来：
\[
X=\begin{bmatrix}x_1^\top\\x_2^\top\\\vdots\\x_n^\top\end{bmatrix}\in\R^{n\times d},\qquad
y=(y_1,\ldots,y_n)^\top\in\R^n.
\]
因此\key{一行是一个样本，一列是一个特征}。$X^\top$ 是转置，即行列互换。$x^\top w=\sum_jx_jw_j$ 是点积，也就是线性模型的预测核心。$\|v\|_2^2=\sum_i v_i^2$ 是平方和；$\|v\|_1=\sum_i|v_i|$ 是绝对值之和。

假设 3 套房子的面积与房龄分别是 $(50,2),(80,8),(60,5)$，则 $X$ 形状为 $3\times2$。如果权重 $w=(2,-1)^\top$，则 $Xw$ 形状为 $3\times1$，第一套的值为 $50\cdot2+2\cdot(-1)=98$。先检查形状，往往比先记矩阵公式更有效。

\section{平均、方差和条件期望}
均值是数据的平均水平。方差衡量围绕均值的平方波动。$\E[Y\mid X=x]$ 读作“已知输入为 $x$ 时，$Y$ 的平均值”。它不同于 $\E[Y]$：后者不看输入，给所有人同一个预测。在平方损失下，条件期望是理论上最好的预测函数；在绝对损失下，对应的是条件中位数。

\section{导数与梯度的直观理解}
导数 $dJ/dw$ 描述 $w$ 小幅增加时 $J$ 如何变化。负号方向通常让函数下降。多个参数时，各方向导数组成梯度 $\nabla J$。梯度下降更新 $w\leftarrow w-\eta\nabla J(w)$，$\eta>0$ 是学习率。学习率太大可能越过谷底，太小则走得慢。

\section{Python 最小工具箱}
在下面的代码中，列表如 \texttt{[1, 2, 3]} 保存一组值；\texttt{for} 重复执行；\texttt{sum} 求和；\texttt{len} 得到个数；\texttt{print} 显示结果。\texttt{import numpy as np} 引入 NumPy 并给它短名 \texttt{np}。\texttt{X.shape} 告诉你矩阵形状。\texttt{model.fit(X, y)} 用数据训练；\texttt{model.predict(X\_new)} 预测新输入。两次调用的输入必须有相同的特征列顺序。

\begin{lstlisting}
prices = [200, 220, 250]
mean_price = sum(prices) / len(prices)
for price in prices:
    print(price - mean_price)
\end{lstlisting}
先在 Python 解释器运行这段不依赖第三方库的代码。后文完整示例要安装 NumPy 与 scikit-learn；可在自己的虚拟环境中执行命令 \texttt{python -m pip install numpy scikit-learn}。

\chapter{从数据到可检验的学习问题}
\section{数据不是现实本身}
监督学习数据集记作 $D=\{(x_i,y_i)\}_{i=1}^{n}$。大写 $X,Y$ 表示随机变量，小写 $x_i,y_i$ 表示已经观察到的值。常用的独立同分布（i.i.d.）假设意味着样本由相同机制产生且彼此独立。它便于分析，却可能被时间趋势、同一家庭的关联样本、城市差异或采集方式破坏。

假如只采集市中心已成交房屋，再把模型用于远郊挂牌房屋，就同时改变了地理范围与价格定义。训练数据覆盖范围、缺失群体、测量误差和标签生成方式都需要记录。

\section{训练误差与未来误差}
单样本损失为 $\ell(y,f(x))$。训练集上的平均损失叫\key{经验风险}：
\[
\widehat R_n(f)=\frac1n\sum_{i=1}^{n}\ell(y_i,f(x_i)).
\]
对未来数据生成分布 $P$ 的期望叫\key{总体风险}：
\[
R(f)=\E_{(X,Y)\sim P}[\ell(Y,f(X))].
\]
训练时能直接计算前者，真正希望降低的是后者。训练集上几乎零误差可以来自记住数据，而不是学到可迁移规律。

\section{训练、验证、测试三种用途}
训练集拟合 $w,b$ 等\key{参数}；验证集选 $K$、多项式阶数、正则化强度等\key{超参数}；测试集只用于最终的一次评估。若反复看测试成绩并据此改模型，测试集就被“用来训练决策”了，分数会偏乐观。小数据可在训练部分使用交叉验证，但最后仍应保留独立测试证据。时间序列应按时间划分；同一个人的多条记录通常要按人分组，避免泄漏。

\section{欠拟合、过拟合与偏差—方差}
模型太简单，训练与测试都差，叫欠拟合；训练很好而新数据很差，叫过拟合。把同一算法在不同训练集上重复拟合：平均预测与真实条件均值的差叫偏差，预测随训练集变动的幅度叫方差。在固定输入 $x$、平方损失和适当独立性条件下，
\[
\begin{aligned}
\E_{D,Y\mid x}\bigl[(Y-\hat f_D(x))^2\bigr]
&=\underbrace{\bigl(\E_D[\hat f_D(x)]-m(x)\bigr)^2}_{\text{偏差平方}}\\
&\quad+\underbrace{\Var_D(\hat f_D(x))}_{\text{模型方差}}
+\underbrace{\Var(Y\mid X=x)}_{\text{不可约噪声}},\\
m(x)&=\E[Y\mid X=x].
\end{aligned}
\]
增大 KNN 的 $K$、降低多项式阶数、加大正则化强度常会降低方差、增加偏差。方向通常如此，具体最优值仍要靠验证集确定。

\begin{remember}预测关系与因果关系不同。面积系数为正，只说明在模型与样本条件下面积和价格有关；“把某套房子的面积增加一平方米会使售价增加多少”需要更强的因果设计。\end{remember}

\chapter{回归：连续数值预测}
\section{什么是回归}
回归预测一个连续数值：房价、耗电量、温度或交付时长。分类预测离散类别或类别概率，如“是否垃圾邮件”。预测是估计未知数值；决策还要结合价格、成本、风险与约束。即便模型预测某房售价为 200 万元，也不能单凭此数决定是否买入。

\section{回归函数与不可约噪声}
定义 $m(x)=\E[Y\mid X=x]$ 和 $\varepsilon=Y-m(X)$，于是
\[
Y=m(X)+\varepsilon,\qquad \E[\varepsilon\mid X]=0.
\]
这个写法不要求 $m$ 线性。即使完全知道 $m$，同样的输入仍可能出现不同输出：交易谈判、测量误差、未记录特征都会引起波动。因此模型误差无法无限降为零。

\section{损失与评价指标}
设误差 $e_i=y_i-\hat y_i$。常用指标为
\[
\MAE=\frac1n\sum_i|e_i|,\quad
\MSE=\frac1n\sum_i e_i^2,\quad
\RMSE=\sqrt{\MSE},\quad
R^2=1-\frac{\sum_i e_i^2}{\sum_i(y_i-\bar y)^2}.
\]
MAE 与 RMSE 的单位和 $y$ 一样，MSE 的单位是 $y$ 单位的平方。大误差在平方损失中权重更高。$R^2=0$ 表示在\emph{同一评估数据}上与用该数据的均值作常数预测一样，$R^2<0$ 表示更差；若分母为零，通常不能按此式解释。MAPE 对真值除法，$y_i=0$ 时无定义，接近零时会异常放大。

\textbf{手算：}$y=(3,5,2)$，$\hat y=(2,5,4)$，绝对误差是 $(1,0,2)$，所以 $\MAE=1$；平方误差和为 $5$，所以 $\RMSE=\sqrt{5/3}\approx1.291$。$\bar y=10/3$，总平方和为 $14/3$，因而 $R^2=1-5/(14/3)=-1/14\approx-0.071$。这里的负分说明这三个预测甚至不如常数基线。

选择指标要依据任务代价：预测高估和低估若有不同后果，单一对称指标可能不足。报告指标时还要给样本量、划分方式与目标单位。

\chapter{K 近邻回归：从“像谁”开始}
\section{算法与手算}
对新输入 $x$，计算它到每个训练输入 $x_i$ 的距离，选最近的 $K$ 个，平均它们的 $y_i$：
\[
\hat f_K(x)=\frac1K\sum_{i\in N_K(x)}y_i.
\]
例如有下面 5 个“面积、离中心距离、月租”样本。查询 $(70,2.5)$，若\emph{暂不标准化}并用欧氏距离，最近三点是第 2、3、1 行，预测 $(2500+2800+3000)/3\approx2766.7$ 元。这个数只是演示距离计算，不是实际租金估计。
\begin{center}\begin{tabular}{ccccc}\toprule
样本&面积（平方米）&距离（公里）&月租（元）&到查询点的距离\\\midrule
1&50&1&3000&20.06\\
2&80&3&2500&10.01\\
3&60&2&2800&10.01\\
4&100&5&2000&30.10\\
5&40&1.5&3200&30.02\\\bottomrule\end{tabular}\end{center}

欧氏距离是 $d(x,z)=\sqrt{\sum_j(x_j-z_j)^2}$；曼哈顿距离是 $\sum_j|x_j-z_j|$。距离定义和 $K$ 都是要选择的超参数。小 $K$ 易受单个异常样本影响，大 $K$ 易把局部差异平均掉。$K=n$ 时预测恒为训练目标均值。

\section{为什么必须关注特征尺度}
面积可能以几十平方米变化，距离只在几公里变化。直接平方相加，面积往往主导“近”。标准化每列为
\[
z_{ij}=\frac{x_{ij}-\mu_j}{s_j},
\]
其中 $\mu_j,s_j$ \key{只从训练集计算}；验证、测试和新数据都复用这些数。若在划分前对全部数据标准化，测试数据的信息已进入训练流程，属于数据泄漏。若某列标准差为零，应移除或使用可处理常数列的预处理器。

KNN 几乎不需训练，但预测时需查找历史样本；它在训练范围外的外推能力很弱。高维时许多点的距离会变得难区分，需要特征选择、降维或另一类模型。

\chapter{线性回归：从直线到矩阵}
\section{模型、截距、残差}
一个特征时 $\hat y=wx+b$；多个特征时 $\hat y=b+\sum_{j=1}^{d}w_jx_j$。$w_j$ 表示在模型中保持其他特征不变时，该特征增加一个单位，预测值怎样变化；它本身不是因果效应。残差是 $e_i=y_i-\hat y_i$。把 $X$ 左边加一列全 1 得到 $\widetilde X$，把 $b$ 放进参数向量 $\beta$，就有 $\hat y=\widetilde X\beta$。

\section{最小二乘及其解}
令目标函数
\[
J(\beta)=\frac{1}{2n}\|\widetilde X\beta-y\|_2^2.
\]
前面的 $1/2$ 让求导后的 $2$ 抵消，不改变最优解。梯度为
\[
\nabla J(\beta)=\frac1n\widetilde X^\top(\widetilde X\beta-y).
\]
令梯度为零，得到\key{正规方程}
\[
\widetilde X^\top\widetilde X\hat\beta=\widetilde X^\top y.
\]
若 $\widetilde X$ 列满秩，$\hat\beta=(\widetilde X^\top\widetilde X)^{-1}\widetilde X^\top y$。目标函数是凸二次函数，所以满足正规方程的解是全局最优；列不满秩时可存在多个系数解，但最小二乘预测可由伪逆给出。

\textbf{务必核算的三点例子：}课件给了 $(1,3),(2,5),(4,4)$，但课件该页印出的拟合直线和残差数值与这三个点\emph{不一致}。直接用均值法：$\bar x=7/3$，$\bar y=4$，
\[
\hat w=\frac{\sum_i(x_i-\bar x)(y_i-\bar y)}{\sum_i(x_i-\bar x)^2}=\frac{1}{14/3}=\frac3{14},\qquad
\hat b=\bar y-\hat w\bar x=\frac72.
\]
所以正确的直线是 $\hat y=\frac72+\frac3{14}x$。三点预测为 $26/7,55/14,61/14$，残差为 $-5/7,15/14,-5/14$；残差之和为 $0$。你可以把三项平方相加，再尝试轻微改变斜率，看平方和是否增大。

\section{几何与数值计算}
$\hat y=\widetilde X\hat\beta$ 是把 $y$ 投影到 $\widetilde X$ 的列空间；残差向量与该列空间正交，即 $\widetilde X^\top(y-\hat y)=0$。几何上这解释了正规方程。显式计算 $(X^\top X)^{-1}$ 在数值上可能放大误差，尤其是特征高度相关或尺度差异很大时。实际软件通常使用 QR、SVD 或其他稳定求解方法；做手算时写逆矩阵公式，写程序时优先使用稳定求解器。标准化能改善尺度问题，但不能自动消除真正的共线性。

\section{统计假设与解释边界}
平方损失最小化\emph{本身}不要求噪声正态。若另加 $Y=X\beta+\varepsilon$、$\E[\varepsilon\mid X]=0$、各样本误差不相关且同方差等条件，便可讨论 OLS 的无偏性与 Gauss--Markov 定理；进一步假设正态通常方便建立精确的有限样本检验。异方差或相关误差会影响常规标准误，遗漏变量可能让系数解释失效。对零基础读者，先会检查残差、预测区间与数据来源，再学习显著性检验。

\chapter{梯度下降：不会一步算出时怎样找答案}
对 $J(\beta)=\|X\beta-y\|_2^2/(2n)$，梯度下降是
\[
\beta^{(t+1)}=\beta^{(t)}-\eta\frac1nX^\top(X\beta^{(t)}-y).
\]
每一步先用当前参数预测，再算残差，再把梯度反方向加到参数上。批量梯度下降每步用全部样本；随机梯度下降每步用一个样本；小批量版本介于两者之间。标准化常让不同方向的坡度更接近，便于设置学习率。对这个二次目标，若 $\eta$ 太大可能发散；若太小会收敛得很慢。常见的检查方式是记录每轮损失，确认总体下降，并同时检查验证集表现。

\textbf{极简演练：}固定截距为零，只有一个样本 $(x,y)=(1,2)$，$J(w)=\frac12(w-2)^2$，梯度 $w-2$。从 $w_0=0$，取 $\eta=0.5$，得到 $w_1=1$、$w_2=1.5$、$w_3=1.75$，逐步靠近 2。这个例子让“负梯度方向”有了可算的含义。

\chapter{复杂模型、正则化与概率解释}
\section{多项式仍可用线性回归训练}
把原输入 $x$ 变为 $\phi(x)=(1,x,x^2,\ldots,x^k)$，模型 $\hat y=\beta^\top\phi(x)$ 对 $x$ 是曲线，但对参数 $\beta$ 仍线性，因此最小二乘推导不变。阶数升高会增加表达能力，也可能放大噪声和数值不稳定。强相关的高次幂使矩阵接近奇异；可考虑标准化、较低阶数、合适基函数和正则化。

\section{Ridge、Lasso、Elastic Net}
为统一推导，以下把\emph{不惩罚截距}的特征矩阵记为 $X$，目标值已按需要中心化，定义
\[
J_{\rm ridge}(w)=\frac12\|Xw-y\|_2^2+\frac\lambda2\|w\|_2^2,
\qquad
J_{\rm lasso}(w)=\frac12\|Xw-y\|_2^2+\lambda\|w\|_1.
\]
Ridge 的方程是 $(X^\top X+\lambda I)\hat w=X^\top y$。$\lambda>0$ 时矩阵正定，能稳定缩小系数，通常不会把某个系数\emph{精确}压成零。Lasso 使用绝对值惩罚，可以让部分系数恰为零；没有类似 Ridge 的通用矩阵闭式解。Elastic Net 同时包含两种惩罚，适合需要稀疏性且特征成组相关的场景。

\begin{center}\begin{tabular}{p{2.6cm}p{3.7cm}p{3.5cm}p{3.6cm}}\toprule
方法&系数效果&相关特征&常见用途\\\midrule
Ridge&整体缩小，通常不精确归零&常共同保留&多数特征有用，系数不稳\\
Lasso&部分系数精确归零&可能择其一，选择不稳&希望得到较少特征\\
Elastic Net&缩小且可归零&更倾向共同保留&特征分组相关\\\bottomrule\end{tabular}\end{center}

\textbf{量纲很重要：}正则化前通常要标准化特征，否则同样的惩罚对不同单位不公平。截距通常不罚。$\lambda$ 的数值依赖前面损失是否除以 $n$、是否写 $1/2$；课件、PDF 和软件包中的 \texttt{alpha} 不能看到数字就直接等同。应先对齐目标函数，再比较数值。

\section{为什么高斯假设会导出平方损失}
如果假设各样本在给定 $x_i$ 下独立，且 $Y_i\mid x_i\sim\mathcal N(x_i^\top w,\sigma^2)$，则负对数似然中与 $w$ 有关的部分是
\[
-\log L(w)=\frac{1}{2\sigma^2}\sum_i(y_i-x_i^\top w)^2+\text{与 }w\text{ 无关的常数}.
\]
所以最大似然估计 $w$ 与最小二乘同解。这个推导是\emph{一种解释}，不是使用最小二乘必须满足的先决条件。若噪声重尾，平方误差可能过度重视离群点，需要检查残差并考虑其他损失。

\section{进阶选读：软阈值与 ISTA}
对一个数 $z$，软阈值算子是
\[
S_\tau(z)=\operatorname{sign}(z)\max(|z|-\tau,0).
\]
例如 $S_2(5)=3$，$S_2(1)=0$，$S_2(-4)=-2$。若 $g(w)=\frac12\|Xw-y\|_2^2$，Lasso 可用迭代软阈值算法（ISTA）：
\[
u^{(t)}=w^{(t)}-\eta X^\top(Xw^{(t)}-y),\qquad
w_j^{(t+1)}=S_{\eta\lambda}(u_j^{(t)}).
\]
先沿光滑的平方损失下降，再让每个坐标向零收缩。步长可取 $0<\eta\le 1/\|X\|_2^2$，这里 $\|X\|_2$ 是最大奇异值。实践中还要规定停止条件，例如相邻两步参数变化足够小；若目标值反而升高，先检查步长和公式约定。

\chapter{完整的回归工作流}
\begin{enumerate}
\item \textbf{定义问题：}要预测什么，何时预测，预测时确实可取得哪些特征，误差代价如何。
\item \textbf{核对数据：}单位、缺失、异常、样本覆盖范围和标签生成过程。
\item \textbf{先划分：}保留独立测试集；有时间或群组结构时相应划分。
\item \textbf{建立基线：}先试“总预测训练均值”，再试简单线性模型，最后才增加复杂度。
\item \textbf{预处理与建模：}仅在训练折拟合标准化或填补器；用流水线连接预处理与模型。
\item \textbf{验证选择：}比较 KNN、线性回归、Ridge 等模型及其超参数，关注方差和样本量。
\item \textbf{最终测试：}报告 MAE、RMSE、$R^2$ 和划分方式；检查残差是否随预测值、时间或群组系统变化。
\item \textbf{部署监控：}记录输入分布、缺失率与误差；出现分布漂移时重新评估。
\end{enumerate}
\textbf{两类常见失败：}数据泄漏指特征包含预测时拿不到的信息，例如用按成交价计算的中介费预测成交价；分布漂移指训练环境和实际使用环境不同，例如仅用市区房屋训练却预测远郊。两者都可能让离线分数失去实际意义。

\chapter{跟着敲两段 Python}
\section{不用第三方包，手算误差}
保存为 \texttt{metrics.py} 后运行 \texttt{python metrics.py}；也可逐行粘贴到 Python 解释器。代码中每一步都对应第 4 章公式。
\begin{lstlisting}
y_true = [3, 5, 2]
y_pred = [2, 5, 4]
n = len(y_true)
errors = [a - b for a, b in zip(y_true, y_pred)]
mae = sum(abs(e) for e in errors) / n
mse = sum(e * e for e in errors) / n
rmse = mse ** 0.5
y_mean = sum(y_true) / n
tss = sum((y - y_mean) ** 2 for y in y_true)
r2 = 1 - sum(e * e for e in errors) / tss
print(mae, mse, rmse, r2)
# Expected: 1.0, 1.666..., 1.29099..., -0.071428...
\end{lstlisting}
注意 \texttt{zip} 会把两组列表按位置配对，\texttt{** 0.5} 是开平方。若两组长度不等，\texttt{zip} 会静默截短，正式代码应先检查长度相等。

\section{可重复的模拟数据练习}
以下代码生成一批\emph{模拟}房价，只用来观察完整流程；输出数值不代表真实房地产结论。安装依赖后在同一脚本中运行。\texttt{random\_state} 固定随机种子，便于复现；\texttt{Pipeline} 让标准化只在各训练折学习。$R^2$ 是评估指标，不是每一步的训练损失。
\begin{lstlisting}
import numpy as np
from sklearn.model_selection import train_test_split, GridSearchCV
from sklearn.pipeline import make_pipeline
from sklearn.preprocessing import StandardScaler
from sklearn.linear_model import LinearRegression, Ridge
from sklearn.neighbors import KNeighborsRegressor
from sklearn.metrics import mean_absolute_error, mean_squared_error, r2_score

rng = np.random.default_rng(42)
area = rng.uniform(45, 140, size=120)
age = rng.uniform(0, 25, size=120)
distance = rng.uniform(1, 15, size=120)
X = np.column_stack([area, age, distance])
noise = rng.normal(0, 12, size=120)
y = 30 + 2.8 * area - 1.5 * age - 4.0 * distance + noise

X_train, X_test, y_train, y_test = train_test_split(
    X, y, test_size=0.2, random_state=42
)
models = {
    "mean": None,
    "linear": LinearRegression(),
    "ridge": make_pipeline(StandardScaler(), Ridge(alpha=10.0)),
    "knn": make_pipeline(StandardScaler(), KNeighborsRegressor(n_neighbors=5)),
}
for name, model in models.items():
    if model is None:
        pred = np.full(len(y_test), y_train.mean())
    else:
        model.fit(X_train, y_train)
        pred = model.predict(X_test)
    mae = mean_absolute_error(y_test, pred)
    rmse = np.sqrt(mean_squared_error(y_test, pred))
    r2 = r2_score(y_test, pred)
    print(name, "MAE", round(mae, 2), "RMSE", round(rmse, 2), "R2", round(r2, 3))

# Only train folds are used during parameter selection.
search = GridSearchCV(
    make_pipeline(StandardScaler(), KNeighborsRegressor()),
    {"kneighborsregressor__n_neighbors": [3, 5, 9, 15]},
    scoring="neg_mean_absolute_error", cv=5
)
search.fit(X_train, y_train)
print("selected K:", search.best_params_)
print("test MAE:", mean_absolute_error(y_test, search.predict(X_test)))
\end{lstlisting}
把 \texttt{noise} 的标准差从 12 改成 40，再观察各方法；或者把真实关系中的 \texttt{area} 改为平方项，看看线性模型何时开始欠拟合。\textbf{不要}根据最后一行测试分数反复挑 $K$，那会把测试集变成验证集。这里比较多个模型的测试输出仅为教学演示；正式评估应在训练数据内部完成所有选择后，只评估一次最终模型。

\chapter{自测题与参考答案}
\section*{题目}
\begin{enumerate}
\item 给定任务“预测未来一天用电量”，写出 $T,P,E$，并说明输出是回归还是分类。
\item $X$ 有 100 行、3 列，$w$ 有 3 个元素。$Xw$ 的形状是什么？若加截距应怎样处理？
\item 训练 MAE 为 1、测试 MAE 为 20，可能发生了什么？至少说出两个检查方向。
\item 为什么计算 KNN 距离之前要标准化？标准化参数应从哪部分数据计算？
\item 在 $(1,3),(2,5),(4,4)$ 上，最小二乘直线的斜率与截距是多少？
\item 如果 $R^2=-0.2$，这在当前评估集上意味着什么？
\item Ridge 和 Lasso 哪个更容易得到精确的零系数？为什么？
\item 从 $w_0=0$ 出发，对 $J(w)=\frac12(w-2)^2$ 取 $\eta=0.5$，算两步梯度下降。
\item 用“未来才知道的维修费”预测当前房价有什么问题？
\end{enumerate}
\section*{答案与理由}
\begin{enumerate}
\item $T$ 是预测下一天的连续用电量，$P$ 可用 MAE（单位千瓦时），$E$ 是历史用电、天气和日历等已知记录；这是回归。
\item $Xw$ 有 100 个预测值，形状通常写作 $100\times1$ 或 NumPy 的 \texttt{(100,)}。加截距可用 $Xw+b$，也可给 $X$ 加全 1 列并把 $b$ 放入参数。
\item 可能过拟合、泄漏、划分方式不当或分布变化。先核对训练/测试划分、预处理是否在训练集内拟合，再检查两边目标值与特征范围。
\item 不同量纲会不公平地支配距离；均值与标准差只从训练数据计算，验证/测试复用。
\item $\hat w=3/14$，$\hat b=7/2$。可代入三点求残差并确认残差和为零。
\item 当前模型的平方误差和比“总预测该评估集均值”的基线高 20\%。不能单凭这个数字判断原因。
\item Lasso。绝对值惩罚的软阈值会把小系数精确压到零；Ridge 常把系数连续缩小。
\item $w_1=1$，$w_2=1.5$；目标值也从 2 降到 $0.5$ 再降到 $0.125$。
\item 预测时拿不到维修费，属于时间上的数据泄漏；离线分数可能虚高。
\end{enumerate}

\chapter{一页公式索引与术语表}
\begin{center}\begin{longtable}{p{4cm}p{9.6cm}}\toprule
概念&公式或含义\\\midrule\endfirsthead\toprule 概念&公式或含义\\\midrule\endhead
经验风险&$\widehat R_n(f)=n^{-1}\sum_i\ell(y_i,f(x_i))$\\
平方损失下的最佳回归函数&$m(x)=\E[Y\mid X=x]$\\
线性模型&$\hat y=Xw+b$，其中矩阵按行放样本\\
最小二乘&$\min_w\|Xw-y\|_2^2$\\
正规方程&$X^\top X\hat w=X^\top y$；含截距时使用增广矩阵\\
梯度下降&$w\leftarrow w-\eta\nabla J(w)$\\
Ridge&平方损失加 $\lambda\|w\|_2^2$；系数收缩\\
Lasso&平方损失加 $\lambda\|w\|_1$；可产生零系数\\
MAE / RMSE&平均绝对误差 / 均方根误差，均与目标同单位\\
$R^2$&$1-\mathrm{RSS}/\mathrm{TSS}$，可为负\\
过拟合&训练集误差低而新数据误差高\\
数据泄漏&预测时不可获得的信息进入训练特征或预处理\\
分布漂移&实际输入或目标机制偏离训练数据\\
\bottomrule\end{longtable}\end{center}

\chapter{资料对应与继续学习}
\section{这份讲义怎样对应原始资料}
\begin{itemize}
\item \textbf{《机器学习简介 · 第一讲 (1).pptx》：}第 1 章对应案例、TPE 定义与学习范式；第 3 章对应数据、经验风险、泛化及数据划分；第 10 章对应系统工作流。其大模型、视觉、自动驾驶、机器人部分在第 1 章作概念地图，避免在回归入门前展开过多细节。
\item \textbf{《机器学习导论 · 第 2 讲：回归.pptx》：}第 4--10 章对应回归目标、KNN、线性回归、梯度下降、多项式、正则化及工作流。第 6 章重新核算了该课件三点最小二乘例子。
\item \textbf{《2-regression.pdf》：}第 3--4 章对应条件期望、训练/测试风险及偏差—方差；第 6--9 章对应 OLS、梯度下降、Ridge/Lasso 和 ISTA 的推导。本文将 ISTA 标记为进阶选读，便于先建立直觉。
\end{itemize}
\section{下一步学什么}
掌握这两讲后，下一站通常是分类：先理解逻辑回归输出概率，再学交叉熵与精确率、召回率；随后学习交叉验证、决策树、聚类和 PCA。复习时反复使用本讲的共同问题：\emph{数据从哪里来？模型允许什么规律？优化什么？如何验证？在哪些情境会失败？}

\section{软件资料}
本讲 Python 示例按 scikit-learn 官方 API 编写；安装包和版本会变化，运行前可查官方文档：\url{https://scikit-learn.org/stable/}。\texttt{train\_test\_split}、\texttt{Pipeline} 与 \texttt{GridSearchCV} 的详细参数见该站 API 文档。若在虚拟环境中出现导入错误，先检查当前 \texttt{python} 与安装包是否属于同一个环境。

\end{document}
